% Preamble
\documentclass[11pt, aspectratio=169]{beamer}
	% Packages
	\usepackage[english]{babel}
	\usepackage{lipsum}
	\usepackage{mathptmx}
	\usepackage{xcolor}
	% Themes
	\usefonttheme[onlymath]{serif}
	\usetheme[nofirafonts]{focus}
	\renewcommand{\thempfootnote}{\arabic{mpfootnote}}
	% Cover
	\title{Microeconometrics Using Stata}
	\subtitle{Linear panel - data models: basics}
	\author{
		Jhon R. Ordoñez
		\footnote{\label{1} National University of San Cristóbal de Huamanga}
		\footnote{\label{2} Faculty of Economic, Administrative and Accounting Sciences}
		\footnote{\label{3} Professional School of Economics}
	}
	\date{\today}
% Body
\begin{document}
	% Cover
	\begin{frame}
		\maketitle
	\end{frame}
	% Outline
	\begin{frame}[t]
		\frametitle{Outline}
		\tableofcontents
	\end{frame}
	% Section 1
	\section{Exercise 1}
		% Slide 1
		\begin{frame}
			\frametitle{Exercise 1}
				\begin{block}{Exercise 1}
					For the data of \textcolor{red!50}{\textbf{section 8.3}}, use \textcolor{blue}{\textbf{xtsum}} 
					to describe the variation in \textcolor{blue}{\textbf{occ}},
					\textcolor{blue}{\textbf{smsa}}, \textcolor{blue}{\textbf{ind}}, \textcolor{blue}{\textbf{ms}},
					\textcolor{blue}{\textbf{union}}, \textcolor{blue}{\textbf{fem}},
					and \textcolor{blue}{\textbf{blk}}. Which of these variables are time
					invariant? Use \textcolor{blue}{\textbf{xttab}} and \textcolor{blue}{\textbf{xttrans}} to provide interpretations of how
					\textcolor{blue}{\textbf{occ}} changes for individuals over the seven years. Provide a time-series
					plot of \textcolor{blue}{\textbf{exp}} for the first $10$ observations, and provide interpretation.
					Provide a scatterplot of \textcolor{blue}{\textbf{lwage}} against \textcolor{blue}{\textbf{ed}}. Is this plot showing within
					variation, between variation, or both?
				\end{block}
		\end{frame}
	% Section 2
	\section{Exercise 2}
		% Slide 1
		\begin{frame}
			\frametitle{Exercise 2}
			\begin{block}{Exercise 2}
				For the data of \textcolor{red!50}{\textbf{section 8.3}}, manually obtain the three standard
				deviations of \textcolor{blue}{\textbf{lwage}} given by the \textcolor{blue}{\textbf{xtsum}} command. For the overall
				standard deviation, use \textcolor{blue}{\textbf{summarize}}. For the between standard
				deviation, compute \textcolor{blue}{\textbf{by id: egen meanwage = mean(lwage)}}, and apply
				\textcolor{blue}{\textbf{summarize}} to \textcolor{blue}{\textbf{(meanwage-grandmean)}} for \textcolor{blue}{\textbf{t==1}}, where grandmean is
				the grand mean over all observations. For the within standard
				deviation, apply \textcolor{blue}{\textbf{summarize}} to \textcolor{blue}{\textbf{(lwage-meanwage)}}. Compare your
				standard deviations with those from \textcolor{blue}{\textbf{xtsum}}. Does
					\begin{equation}
						s_{O}^{2} \simeq s_{W}^{2} + s_{B}^{2}?
					\end{equation}
			\end{block}
		\end{frame}
	% Section 3
	\section{Exercise 3}
		% Slide 1
		\begin{frame}
			\frametitle{Exercise 3}
			\begin{block}{Exercise 3}
				For the model and data of \textcolor{blue}{\textbf{section 8.4}}, compare PFGLS estimators under
				the following assumptions about the error process: independent,
				exchangeable, \textbf{AR(2)}, and \textbf{MA(6)}. Also, compare the associated
				standard-error estimates obtained by using default standard errors and
				by using cluster-robust standard errors. You will find it easiest if you
				combine results using \textcolor{blue}{\textbf{estimates table}}. What happens if you try to fit
				the model with no structure placed on the error correlations?
			\end{block}
		\end{frame}
	% Section 4
	\section{Exercise 4}
		% Slide 1
		\begin{frame}
			\frametitle{Exercise 4}
			\begin{block}{Exercise 4}
				 For the model and data of \textcolor{red!50}{\textbf{section 8.5}}, obtain the within estimator by
				applying \textcolor{blue}{\textbf{regress}} to \textcolor{red!50}{\textbf{(8.7)}}. Hint: For example, for variable x, type by
				\textcolor{blue}{\textbf{id: egen avex = mean(x)}} followed by \textcolor{blue}{\textbf{summarize}} $x$ and then \textcolor{blue}{\textbf{generate}}
				\textcolor{blue}{\textbf{mdx = x - avex + r(mean)}}. Verify that you get the same estimated
				coefficients as you would with \textcolor{blue}{\textbf{xtreg}}, \textcolor{blue}{\textbf{fe}}.
			\end{block}
		\end{frame}
	% Section 5
	\section{Exercise 5}
		% Slide 1
		\begin{frame}
			\frametitle{Exercise 5}
			\begin{block}{Exercise 5}
				For the model and data of \textcolor{red!50}{\textbf{section 8.6}}, compare the \textbf{RE} estimators that
				were obtained by using xtreg with the \textcolor{blue}{\textbf{re}}, \textcolor{blue}{\textbf{mle}}, and \textcolor{blue}{\textbf{pa}} options and
				\textcolor{blue}{\textbf{xtgee}} with the \textcolor{blue}{\textbf{corr(exchangeable)}} option. Also, compare the
				associated standard-error estimates obtained by using default standard
				errors and by using cluster-robust standard errors. You will find it
				easiest if you combine results using \textcolor{blue}{\textbf{estimates table}}.
			\end{block}
		\end{frame}
	% Section 6
	\section{Exercise 6}
		% Slide 1
		\begin{frame}
			\frametitle{Exercise 6}
			\begin{block}{Exercise 6}
				Consider the RE model output given in \textcolor{red!50}{\textbf{section 8.7}}. Verify that, given
				the estimated values of \textcolor{blue}{\textbf{e$\_$sigma}} and \textcolor{blue}{\textbf{u$\_$sigma}}, application of the
				formulas in that section leads to the estimated values of \textcolor{blue}{\textbf{rho}} and
				\textcolor{blue}{\textbf{theta}}.
			\end{block}
		\end{frame}
	% References
	\begin{frame}[t,allowframebreaks]
		\frametitle{References}
		\bibliography{biblio.bib}
		% Style
		\bibliographystyle{apalike}
		% Authors
		\nocite{cameron2022-I}
		\nocite{cameron2022-II}
	\end{frame}
	% Cover
	\begin{frame}
		\maketitle
	\end{frame}
\end{document}